\section{Conditional succinctly computable separations}\label{sec:prf}
We now distinguish explicit evaluation from explicit deterministic selection of a good table. A seeded family is \emph{succinctly evaluable} if its entries are computable in time polynomial in the example encoding length and the supplied seed length. The result below gives a negligible exceptional set of keys. It does not efficiently recognize good keys or specify an efficiently computable deterministic sequence of them.

\paragraph{Cryptographic assumption.}
Fix a constant $\delta>0$, decreasing it to be less than one if necessary. Assume a polynomial-time bit-output pseudorandom function family $\{f_s\}_{s\in\{0,1\}^\sigma}$, defined on inputs long enough to encode a grid row and point, such that every oracle algorithm running in time at most $2^{\sigma^\delta}$ has negligible distinguishing advantage between $f_s$ for uniform $s$ and a uniform random function. Negligibility is in $\sigma$. We use this assumption directly. The standard one-way-function to pseudorandom-generator and generator to pseudorandom-function constructions supply it from suitably subexponentially secure one-way functions, after polynomial reparameterization and a possible reduction of the positive exponent \citep{HILL99,GGM86}. These general cryptographic reductions are imported; the reduction to sign-rank below is given completely.

Take $n=3b+1$, $N=2^b$, so that the domain size is exactly $2N^3=2^n$. Encode $(a,b',x,y)$ injectively in $O(n)$ bits and pad if needed to the PRF input length. Define $A_s$ to be positive off template incidences and to have sign $2f_s(a,b',x,y)-1$ on incidences. Thus a bit one flips a negative entry to positive.

\begin{samepage}
\begin{lemma}[Small subgrid counting and decision cost]\label{lem:subgrid}
At grid scale $k\ge2$, put $K_k=2k^3$ and
\[
 I_k=2k^4-\left(\frac{k(k+1)}2\right)^2,
 \qquad D_k=\left\lfloor\frac{k}{24\log_2 k}\right\rfloor.
\]
Under independent fair flips,
\[
 \Pr[\sr(A)\le D_k]\le2^{-k^4/2}.
\]
Whether $\sr(A)\le D_k$ can be decided after reading at most $I_k\le2k^4$ oracle bits, in time $2^{O(k^3D_k\log k)}=2^{O(k^4)}$ when $D_k\ge1$. When $D_k=0$, the event is empty and is decided immediately.
\end{lemma}
\end{samepage}
\begin{proof}
For $D_k\ge1$ we have $k\ge32$, and
\[
 \log_2(4eK_k/D_k)\le\log_2(8e k^3)\le4\log_2k.
\]
The strict-pattern bound \eqref{eq:warren} therefore gives at most
\[
 2^{2K_kD_k\cdot4\log_2k}\le2^{(2/3)k^4}
\]
low-rank tables. For $k\ge3$, $(k+1)^2/(4k^2)\le4/9<1/2$, so $I_k\ge(3/2)k^4$. Dividing the preceding count by $2^{I_k}$ gives at most $2^{-5k^4/6}\le2^{-k^4/2}$. If $D_k=0$, no nonempty strict sign matrix has that rank.

For the decision algorithm, query the incidence bits and fill the known off-incidence entries. The condition is equivalent to feasibility of
\[
 A_{ij}\sum_{s=1}^{D_k}U_{is}V_{sj}\ge1\quad(i,j\in[K_k]).
\]
Indeed a strict factorization can be scaled to have all signed scores at least one, since there are finitely many positive signed scores. The converse is immediate. There are $2K_kD_k$ real variables and $K_k^2$ degree-two polynomial inequalities with bounded integer coefficients. The singly exponential existential-real decision theorem has bit complexity $(sd)^{O(v)}L^{O(1)}$ for $s$ polynomials of degree at most $d$, $v$ variables, and coefficient bit length $L$ \citep{Renegar92,Vorobjov21}. Substitution gives $2^{O(k^3D_k\log k)}$. This is a use of the classical decision theorem, not a practical algorithm for the large full table.
\end{proof}

\begin{samepage}
\begin{theorem}[Conditional almost-every-key explicit lower bounds]\label{thm:prf}
Fix a constant $c\ge1$ and put
\[
 \sigma=n^{\lceil(4c+2)/\delta\rceil}.
\]
For every key, entries of $A_s$ and hypotheses in $H_{A_s}$ are evaluable in $\operatorname{poly}_c(n)$ time. Every key has RECTIFY and proper learners with $O_\epsilon(n)$ queries and tolerance $\Omega_\epsilon(1)$. For all sufficiently large admissible $n$, outside a negligible fraction of keys,
\[
 \boxed{\dc(H_{A_s})\ge\frac{n^c}{30\log_2n}.}
\]
Here the learner with $O(n)$ queries may use table-polynomial computation. A succinct polynomial-time implementation with $O(n^2)$ queries is established in Section~\ref{sec:fast}.
\end{theorem}
\end{samepage}
\begin{proof}
The evaluation claim follows from integer arithmetic and one PRF evaluation. The rectangle and proper-learning bounds hold for every monotone flip, without any randomness assumption on its entries.

Take the subgrid scale $k=\lceil c n^c\rceil$, which is at most $N$ for sufficiently large $n$. Restrict both the full row set and point set to $[k]\times[2k^2]$. The template and flip law restricted there are exactly the scale-$k$ construction. An oracle distinguisher reads its incidence bits and accepts precisely when its sign-rank is at most $D_k$. Its time is
\[
 2^{O_c(n^{4c})}\le2^{n^{4c+2}}\le2^{\sigma^\delta}
\]
for sufficiently large $n$. Under a uniform function its acceptance probability is at most $2^{-k^4/2}$ by Lemma~\ref{lem:subgrid}. Under the PRF it is therefore at most $2^{-k^4/2}+\operatorname{negl}(\sigma)$, which is negligible in $n$. Restricting a sign realization cannot increase its rank. Consequently, except on that event, the full dimension is at least $D_k+1>k/(24\log_2k)$.

For large $n$, $k\ge c n^c$ and $\log_2k\le c\log_2n+\log_2(2c)\le(5/4)c\log_2n$. Hence
\[
 D_k+1>\frac{k}{24\log_2k}\ge\frac{n^c}{30\log_2n}.
\]
This proves the displayed constant. The choice $k=n^c$ alone would instead give $n^c/(24c\log_2n)$; its factor $c$ cannot be discarded. Also $N=2^{(n-1)/3}$, rather than $2^{n/3}$, respects the fixed definition $n=\log_2|X|$.
\end{proof}

For the linear implication choose any $c>1$. For any proposed fixed-degree polynomial in $(m,1/\tau,n)$ choose $c$ larger than its degree, holding $\epsilon$ fixed. The lower bound then exceeds that polynomial for all sufficiently large $n$. The polynomial key-length exponent depends on $c$; this statement is a hierarchy of succinctly evaluable families, not one polynomial-key family that beats every polynomial. Choosing instead $k=n^{\log_2n}$ and $\sigma=\lceil k^{6/\delta}\rceil$ makes the same proof give superpolynomial dimension with quasipolynomial seed length and evaluation time. It does not settle polynomial-time deterministic explicit selection in the direction of \citet[Problem~4.1]{HHPTZ22}.

\paragraph{Polynomial advice is a different notion of explicitness.}
The cryptographic theorem concerns a prescribed PRF family and supports the hidden-key result in Section~\ref{sec:secret}; cryptography is not necessary for mere polynomial-advice, almost-every-seed evaluation. Indeed, store the $I_k$ independent subgrid bits explicitly in an advice string and set all other template incidences positive. For $k=\lceil cn^c\rceil$ this advice has polynomial length, each entry is polynomial-time evaluable, and Lemma~\ref{lem:subgrid} gives the same high-probability lower bound. Off the subgrid all hypotheses are positive, and all rows outside it are duplicates of the all-positive row. The effective domain is the $2k^3$ subgrid points plus one common positive point; the effective class has at most $2k^3+1$ rows. The fixed-dimensional template and its no-negative-$2\times2$ property persist after adding this positive point: augment the original factorization by one coordinate, zero on old points, one on the new point, and one on every row. The cap implementation is therefore polynomial in $n$ and uses $O_\epsilon(\log k)$ queries. None of these almost-every-advice statements supplies an efficient deterministic selector of good advice. This distinction prevents treating the conditional PRF statement as a solution of deterministic explicit selection.
